
\subsection{Where the span optimum sits}

The choice $L=100$ is not arbitrary and not a monotone preference for more
history. Two ablations bracket it. Against a 50-step pool the 100-step model
wins 5.1--6.4\% on FD002 and FD004 (Table~\ref{tab:abl}). Extending the input
window to 200 steps instead \emph{loses} 1.88\% on FD003 validation
(folds: $-0.10$, $+0.47$, $+4.09$, $-1.52$, $+6.92$\%), and enlarging the
encoder to $d=96$ with three layers loses 6.2\% on two FD001 folds.

Every training engine is longer than 100 cycles (mean 206 on FD001, 247 on
FD003; 92--98\% exceed 150), so the 200-step arm is not starved of data. The
non-monotonicity is therefore about optimisation, not information: a longer
sequence dilutes self-attention over positions whose marginal value is small,
since degradation concentrates at the end of life. We report the bracket rather
than a single point, because it constrains the claim: the method uses about a
hundred steps, not ``all available history''.
